\section{C\'alculo vectorial}
\label{sec:calculo_vectorial}

Consideramos ahora funciones
$\vb{f}:X\subset\R^n\rightarrow\R^m$, con $n,m\in\N$.

\begin{ejemplo}[Campo escalar]
  \label{ej:campo_escalar}
  $f:X\subset\R^2\rightarrow\R$ con $f(\vb{x})=z$.
\end{ejemplo}

\begin{ejemplo}[Transformaci\'on $\R^2\rightarrow\R^2$]
  \label{ej:transformacion}
  \begin{equation}
    \label{eq:transformacion_polar}
    \vb{T}(\rho,\vph)=
    \begin{pmatrix}
      \rho\cos\vph\\
      \rho\sin\vph
    \end{pmatrix}.
  \end{equation}
\end{ejemplo}

\begin{ejemplo}[Superficie $\R^2\rightarrow\R^3$]
  \label{ej:superficie_esfera}
  \begin{equation}
    \label{eq:esfera_param}
    \vb{M}(\theta,\vph)=
    \begin{pmatrix}
      R\sin\theta\cos\vph\\
      R\sin\theta\sin\vph\\
      R\cos\theta
    \end{pmatrix},
  \end{equation}
  cuyas componentes satisfacen $x^2+y^2+z^2=R^2$.
\end{ejemplo}

\begin{ejemplo}[Curva $\R\rightarrow\R^3$]
  \label{ej:curva_R3}
  \begin{equation}
    \label{eq:curva_R3}
    \vb{C}(t)=
    \begin{pmatrix}
      x(t)\\y(t)\\z(t)
    \end{pmatrix},
    \qquad
    \vb{C}:I\subset\R\rightarrow\R^3 .
  \end{equation}
\end{ejemplo}

\begin{definicion}[Curvas y superficies de nivel]
  \label{def:curvas_nivel}
  Sea $f:X\subset\R^2\rightarrow\R$ escalar. Sus \textbf{curvas de nivel} son
  \begin{equation}
    \label{eq:curva_nivel}
    N_k=\set{(x,y)\in\R^2\such f(x,y)=k},
    \qquad k\in\R .
  \end{equation}
  An\'alogamente, si $f:X\subset\R^3\rightarrow\R$, las \textbf{superficies de
  nivel} son $f(x,y,z)=k$ con $k$ constante.
\end{definicion}

\subsection{L\'imites}
\label{subsec:limites_vectoriales}

\begin{definicion}[L\'imite vectorial]
  \label{def:limite_vectorial}
  Sea $\vb{f}:X\subset\R^n\rightarrow\R^m$. Decimos que
  \begin{equation}
    \label{eq:limite_vectorial}
    \lim_{\vb{x}\rightarrow\vb{a}}\vb{f}(\vb{x})=\vb{L}
  \end{equation}
  si dado $\ep>0$ existe $\de_\ep>0$ tal que si $\vb{x}\in X$ y
  $0<\norm{\vb{x}-\vb{a}}<\de_\ep$, entonces
  $\norm{\vb{f}(\vb{x})-\vb{L}}<\ep$.
\end{definicion}

Intuitivamente: para $\norm{\vb{f}-\vb{L}}$ arbitrariamente peque\~no existe
$\norm{\vb{x}-\vb{a}}$ suficientemente peque\~no. Aqu\'i
\begin{equation}
  \label{eq:componentes_f}
  \vb{f}=
  \begin{pmatrix}
    f_1(x_1,\dots,x_n)\\
    f_2(x_1,\dots,x_n)\\
    \vdots\\
    f_m(x_1,\dots,x_n)
  \end{pmatrix}.
\end{equation}

\begin{ejemplo}[Un l\'imite que no existe]
  \label{ej:limite_direccional}
  Sea $f:\R^2\setminus\set{\vb{0}}\rightarrow\R$ con
  \begin{equation}
    \label{eq:f_direccional}
    f(x,y)=\frac{x^2-y^2}{x^2+y^2}.
  \end{equation}
  En coordenadas polares $x=\rho\cos\vph$, $y=\rho\sin\vph$,
  \begin{equation}
    \label{eq:limite_direccional}
    \begin{aligned}
      \lim_{\vb{x}\rightarrow\vb{0}}f(x,y)
        &=\lim_{\rho\rightarrow 0}
          \frac{\rho^2\cos^2\vph-\rho^2\sin^2\vph}{\rho^2}\\
        &=\cos^2\vph-\sin^2\vph ,
    \end{aligned}
  \end{equation}
  que depende del \'angulo, es decir, de la direcci\'on de acercamiento: vale
  $1$ sobre $y=0$, $-1$ sobre $x=0$ y $0$ sobre $y=\pm x$. Si supusi\'eramos
  $L=1$ y buscamos $\abs{f-1}<\ep$, al aproximarnos por la ruta $x=0$
  obtendr\'iamos $\abs{-1-1}=2$, que no es arbitrariamente peque\~no. El
  l\'imite no existe.
\end{ejemplo}

\begin{teorema}[Unicidad]
  \label{teo:unicidad_limite}
  Si un l\'imite existe, es \'unico: si
  $\lim_{\vb{x}\rightarrow\vb{a}}\vb{f}=\vb{L}$ y
  $\lim_{\vb{x}\rightarrow\vb{a}}\vb{f}=\vb{M}$, entonces $\vb{M}=\vb{L}$.
\end{teorema}

\begin{teorema}[Propiedades algebraicas]
  \label{teo:algebra_limites}
  Sean $\vb{F},\vb{G}:X\subset\R^n\rightarrow\R^m$ vectoriales,
  $f,g:X\subset\R^n\rightarrow\R$ escalares y $k\in\R$. Si
  $\lim_{\vb{x}\rightarrow\vb{a}}\vb{F}=\vb{L}$ y
  $\lim_{\vb{x}\rightarrow\vb{a}}\vb{G}=\vb{M}$, entonces
  \begin{align}
    \label{eq:limite_suma}
    \lim_{\vb{x}\rightarrow\vb{a}}(\vb{F}+\vb{G})&=\vb{L}+\vb{M},\\
    \label{eq:limite_escalar}
    \lim_{\vb{x}\rightarrow\vb{a}}k\vb{F}&=k\vb{L}.
  \end{align}
  Si adem\'as $\lim_{\vb{x}\rightarrow\vb{a}}f=L$ y
  $\lim_{\vb{x}\rightarrow\vb{a}}g=M$, entonces
  $\lim_{\vb{x}\rightarrow\vb{a}}(fg)=LM$, y si $M\neq 0$ y $g$ no se anula
  cerca de $\vb{a}$,
  \begin{equation}
    \label{eq:limite_cociente}
    \lim_{\vb{x}\rightarrow\vb{a}}\frac{f}{g}=\frac{L}{M}.
  \end{equation}
\end{teorema}

\begin{teorema}[L\'imites por componentes]
  \label{teo:limites_componentes}
  Sea $\vb{f}:X\subset\R^n\rightarrow\R^m$. Entonces
  $\lim_{\vb{x}\rightarrow\vb{a}}\vb{f}(\vb{x})=\vb{L}$ con
  $\vb{L}=(L_1,\dots,L_m)^\tp$ si y s\'olo si
  $\lim_{\vb{x}\rightarrow\vb{a}}f_i(\vb{x})=L_i$ para cada $i$.
\end{teorema}

\begin{ejemplo}
  \label{ej:limite_existe}
  Sea $f:\R^2\setminus\set{\vb{0}}\rightarrow\R$ con
  $f(x,y)=\qty{x^3+y^5}/\qty{x^2+y^2}$. En polares,
  \begin{equation}
    \label{eq:limite_polares}
    \begin{aligned}
      \lim_{\vb{x}\rightarrow\vb{0}}f(x,y)
        &=\lim_{\rho\rightarrow 0}
          \frac{\rho^3\cos^3\vph+\rho^5\sin^5\vph}{\rho^2}\\
        &=\lim_{\rho\rightarrow 0}
          \qty{\rho\cos^3\vph+\rho^3\sin^5\vph}=0
    \end{aligned}
  \end{equation}
  para todo $\vph$. Aqu\'i el l\'imite s\'i existe, y la cota es uniforme en la
  direcci\'on: dado $\ep>0$ basta $\rho<\ep/2$, pues
  $\abs{\rho\cos^3\vph+\rho^3\sin^5\vph}\le\rho+\rho^3$.
\end{ejemplo}

\subsection{Continuidad}
\label{subsec:continuidad_vectorial}

\begin{definicion}[Continuidad]
  \label{def:continuidad_vectorial}
  Sea $\vb{f}:X\subset\R^n\rightarrow\R^m$ y $\vb{a}\in X$. Entonces $\vb{f}$
  es \textbf{continua} en $\vb{a}$ si $\vb{a}$ es un punto aislado o si
  \begin{equation}
    \label{eq:continuidad_vectorial}
    \lim_{\vb{x}\rightarrow\vb{a}}\vb{f}(\vb{x})=\vb{f}(\vb{a}).
  \end{equation}
  Si esto ocurre en todos los puntos de su dominio, $\vb{f}$ es continua en
  $X$.
\end{definicion}

\subsection{Derivadas parciales}
\label{subsec:derivadas_parciales}

\begin{definicion}[Derivada parcial]
  \label{def:derivada_parcial}
  La derivada parcial de $\vb{f}$ respecto de $x_i$ es la derivada de la
  funci\'on parcial correspondiente:
  \begin{equation}
    \label{eq:derivada_parcial}
    \small
    \pdv{\vb{f}}{x_i}=\lim_{h\rightarrow 0}
      \frac{\vb{f}(x_1,\dots,x_i+h,\dots,x_n)
            -\vb{f}(x_1,\dots,x_i,\dots,x_n)}{h}.
  \end{equation}
\end{definicion}

\begin{ejemplo}
  \label{ej:parcial_escalar}
  Sea $f:\R^2\rightarrow\R$ con $f(x,y)=x^2y+\cos(x+y)$. Entonces
  \begin{equation}
    \label{eq:parcial_escalar}
    \small
    \begin{aligned}
      \pdv{f}{x}
        &=\lim_{h\rightarrow 0}\frac{f(x+h,y)-f(x,y)}{h}\\
        &=\lim_{h\rightarrow 0}\frac{(x+h)^2y-x^2y}{h}\\
        &\quad+\lim_{h\rightarrow 0}
          \frac{\cos(x+h+y)-\cos(x+y)}{h}\\
        &=2xy-\sin(x+y).
    \end{aligned}
  \end{equation}
\end{ejemplo}

\begin{ejemplo}
  \label{ej:parcial_vectorial}
  Sea $\vb{f}:\R^2\rightarrow\R^2$ con
  \begin{equation}
    \label{eq:f_vectorial}
    \vb{f}(x,y)=
    \begin{pmatrix}
      2x^2+y\\
      x^3\sin(xy)
    \end{pmatrix}.
  \end{equation}
  La derivada parcial act\'ua componente a componente,
  \begin{equation}
    \label{eq:parcial_vectorial}
    \pdv{\vb{f}}{y}=
    \begin{pmatrix}
      \pdv{}{y}\qty{2x^2+y}\\[2pt]
      \pdv{}{y}\qty{x^3\sin(xy)}
    \end{pmatrix}
    =
    \begin{pmatrix}
      1\\
      x^4\cos(xy)
    \end{pmatrix}.
  \end{equation}
\end{ejemplo}

\subsection{Diferenciabilidad}
\label{subsec:diferenciabilidad}

Sean $V\subset\R^n$ y $W\subset\R^m$ espacios normados y
$\vb{f}:V\rightarrow W$. La idea es aproximar $\vb{f}$ por una parte lineal
m\'as un residuo que se anula m\'as r\'apido que $\norm{\vb{r}}$:
\begin{equation}
  \label{eq:aproximacion_lineal}
  \vb{f}(\vb{r})=\vb{g}(\vb{r})+\lito(\vb{r}),
\end{equation}
donde la $\lito$ peque\~na se define por
\begin{equation}
  \label{eq:o_pequena}
  \lim_{\vb{r}\rightarrow\vb{0}}
    \frac{\lito(\vb{r})}{\norm{\vb{r}}}=\vb{0}\in\R^m .
\end{equation}

\subsubsection{Analog\'ia en una dimensi\'on}
\label{subsubsec:analogia_1d}

Para $g:\R\rightarrow\R$, decir que $g$ es diferenciable en $x$ equivale a
\begin{equation}
  \label{eq:diferenciable_1d}
  \lim_{h\rightarrow 0}\frac{1}{h}
    \qty{g(x+h)-g(x)-h\dv{}{x}g}=0,
\end{equation}
o, escrito de otro modo,
\begin{equation}
  \label{eq:diferenciable_1d_bis}
  g(x+h)=g(x)+h\dv{}{x}g(x)+\lito(h).
\end{equation}

\begin{definicion}[Diferenciabilidad]
  \label{def:diferenciabilidad}
  $\vb{f}:V\subset\R^n\rightarrow W\subset\R^m$ es \textbf{diferenciable} si
  existe una transformaci\'on lineal $D\vb{f}(\vb{r}):V\rightarrow W$ tal que
  \begin{equation}
    \label{eq:diferenciabilidad}
    \vb{f}(\vb{r}+\vb{u})
      =\vb{f}(\vb{r})+D\vb{f}(\vb{r})[\vb{u}]+\lito(\vb{u}).
  \end{equation}
\end{definicion}

\begin{teorema}
  \label{teo:unicidad_diferencial}
  Si $\vb{f}$ es diferenciable en $\vb{r}\in\R^n$, hay una \'unica
  transformaci\'on lineal $D\vb{f}(\vb{r}):\R^n\rightarrow\R^m$ tal que
  \begin{equation}
    \label{eq:unicidad_diferencial}
    \lim_{\xii\rightarrow 0}
      \abs{\frac{\vb{f}(\vb{r}+\xii\vb{u})-\vb{f}(\vb{r})}
                {\norm{\xii\vb{u}}}-D\vb{f}(\vb{r})[\vb{u}]}=0 .
  \end{equation}
\end{teorema}

De aqu\'i se obtiene la expresi\'on operativa: la \textbf{derivada
direccional}
\begin{equation}
  \label{eq:derivada_direccional}
  \begin{aligned}
    D\vb{f}(\vb{r})[\vb{u}]
      &=\lim_{\xii\rightarrow 0}\frac{1}{\xii\norm{\vb{u}}}
        \qty{\vb{f}(\vb{r}+\xii\vb{u})-\vb{f}(\vb{r})}\\
      &=\frac{1}{\norm{\vb{u}}}
        \eval{\dv{}{\xii}\vb{f}(\vb{r}+\xii\vb{u})}_{\xii=0}.
  \end{aligned}
\end{equation}

\begin{ejemplo}[Caso escalar]
  \label{ej:gradiente_escalar}
  Sea $f:\R^n\rightarrow\R$. Escribiendo
  $\et_i=x_i+\xii u_i$ y aplicando la regla de la cadena,
  \begin{equation}
    \label{eq:derivada_direccional_escalar}
    \small
    \begin{aligned}
      Df(\vb{r})[\vb{u}]
        &=\frac{1}{\norm{\vb{u}}}
          \eval{\dv{}{\xii}f(\et_1,\dots,\et_n)}_{\xii=0}\\
        &=\frac{1}{\norm{\vb{u}}}
          \eval{\sum_{i=1}^{n}\pdv{f}{\et_i}\dv{\et_i}{\xii}}_{\xii=0}.
    \end{aligned}
  \end{equation}
  Como $\lim_{\xii\rightarrow 0}\et_i=x_i$ y
  $\dv{\et_i}{\xii}=u_i$, resulta
  \begin{equation}
    \label{eq:gradiente_producto}
    Df(\vb{r})[\vb{u}]=\frac{1}{\norm{\vb{u}}}
      \sum_{i=1}^{n}\pdv{f}{x_i}u_i
      =\nabla f\cdot\vh{u}.
  \end{equation}
\end{ejemplo}

\subsection{El gradiente}
\label{subsec:gradiente}

Generalizando el ejemplo anterior a $\vb{f}:\R^n\rightarrow\R^m$, la regla de
la cadena da
\begin{equation}
  \label{eq:jacobiano}
  \small
  D\vb{f}(\vb{r})[\vb{u}]=\frac{1}{\norm{\vb{u}}}
  \begin{pmatrix}
    \pdv{f_1}{r_1} & \dots & \pdv{f_1}{r_n}\\
    \pdv{f_2}{r_1} & \dots & \pdv{f_2}{r_n}\\
    \vdots & \ddots & \vdots\\
    \pdv{f_m}{r_1} & \dots & \pdv{f_m}{r_n}
  \end{pmatrix}
  \begin{pmatrix}
    u_1\\u_2\\\vdots\\u_n
  \end{pmatrix},
\end{equation}
es decir
\begin{equation}
  \label{eq:gradiente_matriz}
  D\vb{f}(\vb{r})[\vb{u}]=\grad\vb{f}\;\frac{\vb{u}}{\norm{\vb{u}}}
    =\grad\vb{f}\,\vh{u},
\end{equation}
donde la matriz de \eqref{eq:jacobiano} es el \textbf{gradiente} (o matriz
jacobiana) de $\vb{f}$.

\begin{ejemplo}
  \label{ej:gradiente_matriz}
  Sea $\vb{f}:\R^2\rightarrow\R^3$ con
  \begin{equation}
    \label{eq:f_ejemplo_gradiente}
    \vb{f}(x,y)=
    \begin{pmatrix}
      3x^2-y^3\\
      xy+2\\
      x-y+xy
    \end{pmatrix}
    \imp
    \grad\vb{f}=
    \begin{pmatrix}
      6x & -3y^2\\
      y & x\\
      1+y & -1+x
    \end{pmatrix}.
  \end{equation}
\end{ejemplo}

\begin{teorema}[Direcci\'on de m\'aximo cambio]
  \label{teo:maximo_cambio}
  Sea $f:\R^n\rightarrow\R$ con $\nabla f\neq\vb{0}$. Entonces
  $\nabla f(\vb{r})$ apunta en la direcci\'on en que $f$ tiene mayor cambio.
\end{teorema}

\begin{ejemplo}
  \label{ej:gradiente_silla}
  Sea $f:\R^2\rightarrow\R$ con $f=x^2-y^2$. Su gradiente es
  \begin{equation}
    \label{eq:gradiente_silla}
    \nabla f=
    \begin{pmatrix}
      \pdv{}{x}\qty{x^2-y^2}\\[2pt]
      \pdv{}{y}\qty{x^2-y^2}
    \end{pmatrix}
    =2
    \begin{pmatrix}
      x\\-y
    \end{pmatrix},
  \end{equation}
  y sobre la recta $y=x$ vale $\nabla f=2x(1,-1)^\tp$.
\end{ejemplo}

\subsection{Plano tangente}
\label{subsec:plano_tangente}

\begin{teorema}[El gradiente es normal a la superficie de nivel]
  \label{teo:gradiente_normal}
  Sea $f:\R^3\rightarrow\R$ un mapeo de clase $C^1$ y $S$ la superficie de
  nivel definida por $f(x,y,z)=k$. Si $\vb{v}$ es el vector tangente en
  $\xii=0$ a una curva $\vb{c}(\xii)$ contenida en $S$ con
  $\vb{c}(0)=(x_0,y_0,z_0)$, entonces
  $\eval{\nabla f}_{(x_0,y_0,z_0)}\cdot\vb{v}=0$.
\end{teorema}

\begin{demostracion}
  La curva vive sobre $f(x,y,z)=k$, luego
  $f(x(\xii),y(\xii),z(\xii))=k$ es constante y su derivada se anula. Por la
  regla de la cadena,
  \begin{equation}
    \label{eq:demo_gradiente_normal}
    \begin{aligned}
      \dv{}{\xii}f
        &=\pdv{f}{x}\dv{x}{\xii}+\pdv{f}{y}\dv{y}{\xii}
          +\pdv{f}{z}\dv{z}{\xii}\\
        &=\nabla f\cdot\dv{}{\xii}\vb{c}
         =\nabla f\cdot\vb{v}=0 .
    \end{aligned}
  \end{equation}
\end{demostracion}

\begin{ejemplo}
  \label{ej:esfera_nivel}
  Sea $f(x,y,z)=x^2+y^2+z^2$. La superficie de nivel $f=R^2>0$ es la esfera de
  radio $R$, y
  \begin{equation}
    \label{eq:gradiente_esfera}
    \nabla\qty{x^2+y^2+z^2}=2
    \begin{pmatrix}
      x\\y\\z
    \end{pmatrix},
  \end{equation}
  que en cada punto es radial, es decir, normal a la esfera.
\end{ejemplo}

Si dos curvas $\vb{c}_1$ y $\vb{c}_2$ contenidas en $S$ pasan por
$\vb{r}_0$ y sus tangentes
\begin{equation}
  \label{eq:tangentes_curvas}
  \vb{v}_1=\eval{\dv{}{\xii}\vb{c}_1}_{\vb{r}_0},
  \qquad
  \vb{v}_2=\eval{\dv{}{\xii}\vb{c}_2}_{\vb{r}_0}
\end{equation}
no son colineales, entonces son linealmente independientes y generan el
\textbf{plano tangente} a $S$ en $\vb{r}_0$. Como la ecuaci\'on de un plano de
normal $\vb{N}$ que pasa por $\vb{r}_0$ es
$\vb{N}\cdot(\vb{r}-\vb{r}_0)=0$, y por el
Teorema~\ref{teo:gradiente_normal} la normal es $\nabla f$, el plano tangente
es
\begin{equation}
  \label{eq:plano_tangente}
  \eval{\nabla f}_{\vb{r}_0}\cdot\qty{\vb{r}-\vb{r}_0}=0 .
\end{equation}

\subsection{Teorema de la funci\'on impl\'icita}
\label{subsec:funcion_implicita}

\begin{teorema}[Funci\'on impl\'icita]
  \label{teo:funcion_implicita}
  Sea $F:X\subset\R^n\rightarrow\R$ de clase $C^1$ y sea $\vb{a}$ un punto del
  conjunto de nivel
  $S=\set{\vb{x}\in\R^n\such F(\vb{x})=k}$. Si
  $\pdv{F}{x_i}(\vb{a})\neq 0$, entonces existen una vecindad $U$ de
  $(a_1,\dots,a_{i-1},a_{i+1},\dots,a_n)\in\R^{n-1}$, una vecindad $V$ de
  $a_i\in\R$ y una funci\'on $f:U\rightarrow V$ de clase $C^1$ tales que si
  $(x_1,\dots,x_{i-1},x_{i+1},\dots,x_n)\in U$ y $x_i\in V$ satisfacen
  $F(x_1,\dots,x_n)=k$, entonces
  $x_i=f(x_1,\dots,x_{i-1},x_{i+1},\dots,x_n)$.
\end{teorema}

\begin{ejemplo}
  \label{ej:elipsoide_implicita}
  Sea
  \begin{equation}
    \label{eq:elipsoide_nivel}
    f(x,y,z)=\frac{x^2}{4}+\frac{y^2}{36}+\frac{z^2}{9}=k .
  \end{equation}
  Tomando la superficie de nivel $k=1$ obtenemos un elipsoide. Como
  \begin{equation}
    \label{eq:parcial_elipsoide}
    \pdv{f}{z}=\frac{2z}{9},
  \end{equation}
  en el punto $\vb{a}=(0,0,-3)^\tp$, que pertenece al elipsoide, se tiene
  \begin{equation}
    \label{eq:parcial_elipsoide_a}
    \eval{\pdv{f}{z}}_{\vb{x}=\vb{a}}=-\frac{2}{3}\neq 0 .
  \end{equation}
  Por el teorema de la funci\'on impl\'icita, alrededor de ese punto podemos
  despejar
  \begin{equation}
    \label{eq:despeje_z}
    z=-3\sqrt{1-\frac{x^2}{4}-\frac{y^2}{36}} .
  \end{equation}
\end{ejemplo}

\begin{teorema}[Funci\'on impl\'icita generalizado]
  \label{teo:funcion_implicita_gen}
  Supongamos que $\vb{F}:A\subset\R^{n+m}\rightarrow\R^m$ es de clase $C^1$ y
  que $(\vb{a},\vb{b})\in A$ satisface $\vb{F}(\vb{a},\vb{b})=\vb{k}$. Si el
  determinante
  \begin{equation}
    \label{eq:determinante_implicita}
    \small
    \det
    \begin{pmatrix}
      \pdv{F_1}{y_1}(\vb{a},\vb{b}) & \dots & \pdv{F_1}{y_m}(\vb{a},\vb{b})\\
      \vdots & \ddots & \vdots\\
      \pdv{F_m}{y_1}(\vb{a},\vb{b}) & \dots & \pdv{F_m}{y_m}(\vb{a},\vb{b})
    \end{pmatrix}
    \neq 0,
  \end{equation}
  entonces hay una vecindad $U$ de $\vb{a}\in\R^n$ y una \'unica funci\'on
  $\vb{f}:U\rightarrow\R^m$ de clase $C^1$ tal que
  \begin{equation}
    \label{eq:conclusion_implicita}
    \vb{f}(\vb{a})=\vb{b},
    \qquad
    \vb{F}(\vb{x},\vb{f}(\vb{x}))=\vb{k}
    \quad\forall\vb{x}\in U .
  \end{equation}
\end{teorema}

\subsection{Teorema de la funci\'on inversa}
\label{subsec:funcion_inversa}

\begin{teorema}[Funci\'on inversa]
  \label{teo:funcion_inversa}
  Supongamos que $\vb{f}:\R^n\rightarrow\R^n$ es de clase $C^1$ sobre un
  abierto $A\subset\R^n$ y que en $\vb{a}\in A$ el jacobiano no se anula,
  \begin{equation}
    \label{eq:jacobiano_inversa}
    \abs{\frac{\partial(f_1,f_2,\dots,f_n)}
              {\partial(x_1,x_2,\dots,x_n)}}\neq 0 .
  \end{equation}
  Entonces existe un abierto $U\subset\R^n$ que contiene a $\vb{a}$ tal que
  $\vb{f}$ es uno a uno sobre $U$, el conjunto $V=\vb{f}(U)$ tambi\'en es
  abierto y hay una \'unica funci\'on inversa $\vb{g}:V\rightarrow U$ de
  $\vb{f}$, tambi\'en de clase $C^1$.
\end{teorema}
